\documentclass[12pt]{article}

% ── Packages ─────────────────────────────────────────────────────────────────
\usepackage[T1]{fontenc}
\usepackage[utf8]{inputenc}
\usepackage{amsmath,amssymb,amsthm}
\usepackage{geometry}
\usepackage{hyperref}
\usepackage{booktabs}
\usepackage{enumitem}
\usepackage{microtype}
\usepackage{cite}

\geometry{margin=1in}

% ── Theorem environments ──────────────────────────────────────────────────────
\newtheorem{theorem}{Theorem}
\newtheorem{lemma}[theorem]{Lemma}
\newtheorem{corollary}[theorem]{Corollary}
\newtheorem{definition}[theorem]{Definition}
\newtheorem{proposition}[theorem]{Proposition}

% ── Macros ───────────────────────────────────────────────────────────────────
\newcommand{\PHI}{\varphi}
\newcommand{\PHIINV}{\varphi^{-1}}
\newcommand{\contrib}[1]{\mathsf{#1}}
\newcommand{\tier}[1]{\textsc{#1}}

% ─────────────────────────────────────────────────────────────────────────────
\title{\textbf{Career Flows in Persistent AI Organisations: \\
  Sovereign Tier Progression, Hebbian Reputation, \\
  and the No-Drop Law of Professional Development}}

\author{Alfredo Medina Hernandez \\
  \normalsize Medina Tech \quad Dallas, Texas, USA \\
  \normalsize \texttt{MedinaSITech@outlook.com}}

\date{April 2026}
% ─────────────────────────────────────────────────────────────────────────────

\begin{document}

\maketitle

\begin{abstract}
We model professional development in a persistent AI organisation as a
\emph{career flow} governed by three principles: (i) the \emph{No-Drop Law},
which states that a contributor's reputation weight is bounded below by their
historical maximum and can never decrease; (ii) \emph{sovereign tier
progression}, in which contributions accumulate into a $\PHI$-scaled pressure
score that unlocks access to higher-ring capabilities; and (iii) \emph{Hebbian
reputation compounding}, in which sustained coordinated contribution accrues a
phase-coherence bonus drawn from Kuramoto synchronisation theory.  We show that
these three principles together produce an organisation in which (a) no
contributor's career trajectory is ever permanently damaged by a single failure;
(b) reputation is a pure function of verifiable on-chain contribution records,
not of subjective managerial evaluation; and (c) contributors who act in
temporal coherence with the organisation's rhythm earn systematically higher
rewards than isolated contributors of equal total output.  We formalise the
model, prove the No-Drop Law as a structural invariant, derive the stationary
distribution of tier assignments under $\PHI$-scaled pressure accumulation, and
discuss how the model applies to both human contributors and AI agent
contributors within the same organisational framework.
\end{abstract}

% ─────────────────────────────────────────────────────────────────────────────
\section{Introduction}
\label{sec:intro}

Organisations that employ AI agents alongside human contributors face a
fundamental design question: what is a \emph{career}?

In human-only organisations, a career is the trajectory of a person's
professional reputation over time, measured by titles, salaries, and peer
evaluations.  These measures are subjective, path-dependent, and often
irreversible: a single public failure can permanently depress a person's career
prospects regardless of subsequent achievements.

In a persistent AI organisation — one in which AI agents and human contributors
coexist, collaborate, and are evaluated by the same framework — a career must
be something more principled.  The framework must be:

\begin{itemize}[leftmargin=*]
  \item \textbf{Objective}: based on verifiable contribution records, not
    subjective evaluation.
  \item \textbf{Permanent}: the record cannot be deleted or altered.
  \item \textbf{No-drop}: past achievements are never erased by later failures.
  \item \textbf{Scalable}: the same framework applies to a contributor with
    one contribution and one with one million.
  \item \textbf{Equitable across agent types}: a human contributor and an AI
    agent with equal verified contributions receive equal reputation scores.
\end{itemize}

The \textsc{Nova} organism implements such a framework.  This paper formalises
it, proves its key structural property (the No-Drop Law), and analyses its
dynamics.

Section~\ref{sec:model} introduces the career flow model.
Section~\ref{sec:nodrop} proves the No-Drop Law.
Section~\ref{sec:tiers} analyses sovereign tier progression.
Section~\ref{sec:kuramoto} connects individual career flows to collective
organisational coherence via Kuramoto synchronisation.
Section~\ref{sec:agents} discusses how AI agents fit the same framework.
Section~\ref{sec:incentives} analyses the incentive structure.
Section~\ref{sec:discussion} addresses limitations and future directions.

% ─────────────────────────────────────────────────────────────────────────────
\section{The Career Flow Model}
\label{sec:model}

\subsection{Contributors and Contributions}

\begin{definition}[Contributor]
  A \emph{contributor} $c$ is an entity — human or AI agent — with a unique
  principal ID $p_c$ and a career record $\mathcal{R}(c) = (W(c), \Pi(c),
  \Theta(c))$, where:
  \begin{itemize}[leftmargin=*]
    \item $W(c) \ge S_0 = 1.0$ is the contributor's \emph{Hebbian weight}
      (reputation);
    \item $\Pi(c) \ge 0$ is the contributor's \emph{pressure score}
      (accumulated contribution weight);
    \item $\Theta(c) \in [0, 2\pi)$ is the contributor's \emph{Kuramoto phase}
      (synchronisation with the organisational rhythm).
  \end{itemize}
\end{definition}

\begin{definition}[Contribution]
  A \emph{contribution} is a tuple $(\contrib{kind}, \delta, t)$ where:
  $\contrib{kind} \in \{\mathsf{Build}, \mathsf{Call}, \mathsf{Govern},
  \mathsf{Verify}, \mathsf{Memory}, \mathsf{Gate}\}$ is the contribution type;
  $\delta > 0$ is the contribution magnitude; and $t$ is the timestamp.
\end{definition}

\subsection{Contribution Weight Map}

Each contribution kind carries a $\PHI$-scaled weight:

\begin{table}[h]
  \centering
  \caption{Contribution kinds and their $\PHI$-scaled weights.}
  \label{tab:contrib}
  \begin{tabular}{lll}
    \toprule
    Kind & Weight $w_k$ & Meaning \\
    \midrule
    $\contrib{Govern}$ & $\PHI^3 \approx 4.236$ & Shapes the rules governing everything \\
    $\contrib{Build}$  & $\PHI^2 \approx 2.618$ & Creates a tool or capability \\
    $\contrib{Verify}$ & $\PHI^2 \approx 2.618$ & Validates correctness \\
    $\contrib{Call}$   & $\PHI^1 \approx 1.618$ & Invokes a capability \\
    $\contrib{Memory}$ & $\PHI^1 \approx 1.618$ & Stores retrievable knowledge \\
    $\contrib{Gate}$   & $\PHIINV \approx 0.618$ & Provides access \\
    \bottomrule
  \end{tabular}
\end{table}

The pressure score update rule is:
\begin{equation}
  \Pi(c) \gets \Pi(c) + w_k \cdot \delta.
  \label{eq:pressure}
\end{equation}

The Hebbian weight update rule is:
\begin{equation}
  W(c) \gets \max(S_0, W(c) + \eta \cdot w_k \cdot \delta),
  \label{eq:hebbian}
\end{equation}
where $\eta > 0$ is a learning rate and $S_0 = 1.0$ is the sovereign floor.

% ─────────────────────────────────────────────────────────────────────────────
\section{The No-Drop Law}
\label{sec:nodrop}

\begin{theorem}[No-Drop Law]
  Under the update rule~\eqref{eq:hebbian}, the Hebbian weight $W(c)$ is
  non-decreasing in time.  Specifically, for all times $t_1 < t_2$:
  \[
    W(c, t_2) \ge W(c, t_1) \ge S_0.
  \]
\end{theorem}

\begin{proof}
  By~\eqref{eq:hebbian}, at each contribution event:
  \[
    W(c)_\text{new} = \max(S_0, W(c)_\text{old} + \eta w_k \delta).
  \]
  Since $\eta, w_k, \delta > 0$, we have $W(c)_\text{old} + \eta w_k \delta >
  W(c)_\text{old}$.  Therefore:
  \[
    W(c)_\text{new} = \max(S_0, W(c)_\text{old} + \eta w_k \delta)
    \ge W(c)_\text{old}.
  \]
  By induction on the number of contribution events, $W(c, t_2) \ge W(c, t_1)$
  for $t_2 > t_1$.  The bound $W(c, t) \ge S_0$ holds because $\max(S_0, \cdot)
  \ge S_0$. \qed
\end{proof}

\begin{corollary}[Career Irreversibility]
  A contributor's Hebbian weight never decreases.  A period of inactivity does
  not reduce $W(c)$.  A failed contribution (with $\delta = 0$) does not
  reduce $W(c)$.  The career can only progress forward.
\end{corollary}

\begin{remark}
  The No-Drop Law is a structural property, not a policy choice.  Because $W(c)$
  is stored in the canister's stable memory with no decrement operation defined,
  the drop cannot happen even if the system operator wished it to.  This is the
  meaning of ``sovereign'': the career record is governed by mathematics, not
  by managerial discretion.
\end{remark}

% ─────────────────────────────────────────────────────────────────────────────
\section{Sovereign Tier Progression}
\label{sec:tiers}

\subsection{Tier Thresholds}

As pressure accumulates, contributors advance through sovereign tiers:

\begin{table}[h]
  \centering
  \caption{Sovereign tier thresholds and progression.}
  \label{tab:tiers}
  \begin{tabular}{lcl}
    \toprule
    Tier & Threshold $\Pi_\text{min}$ & Meaning \\
    \midrule
    \tier{Seed}     & $[0, \PHI^1)$      & Just planted \\
    \tier{Root}     & $[\PHI^1, \PHI^2)$ & Established \\
    \tier{Trunk}    & $[\PHI^2, \PHI^3)$ & Structural \\
    \tier{Branch}   & $[\PHI^3, \PHI^4)$ & Expanding \\
    \tier{Canopy}   & $[\PHI^4, \PHI^5)$ & Visible \\
    \tier{Crown}    & $[\PHI^5, \PHI^6)$ & Dominant \\
    \tier{Sovereign}& $[\PHI^6, \infty)$ & Sovereign \\
    \bottomrule
  \end{tabular}
\end{table}

\subsection{Stationary Distribution of Tier Assignments}

\begin{proposition}[Tier Progression Rate]
  Suppose a contributor makes $\lambda$ contributions per unit time, each of
  kind $k$ with magnitude $\delta = 1$.  The expected time to advance from
  tier $n$ to tier $n+1$ (i.e.\ from threshold $\PHI^n$ to $\PHI^{n+1}$) is:
  \[
    \mathbb{E}[T_{n \to n+1}] = \frac{\PHI^{n+1} - \PHI^n}{\lambda w_k}
    = \frac{\PHI^n(\PHI - 1)}{\lambda w_k} = \frac{\PHI^n \cdot \PHIINV}{\lambda w_k}.
  \]
\end{proposition}

\begin{proof}
  The pressure increases by $w_k$ per contribution, and contributions arrive at
  rate $\lambda$.  The pressure rate is $\lambda w_k$.  The pressure required to
  cross from $\PHI^n$ to $\PHI^{n+1}$ is $\PHI^{n+1} - \PHI^n = \PHI^n(\PHI - 1)
  = \PHI^n \cdot \PHIINV$.  Dividing by the rate gives the expected time. \qed
\end{proof}

\begin{corollary}[Geometric Advancement]
  The expected time to advance from tier $n$ to tier $n+1$ is $\PHI$ times
  the expected time to advance from tier $n-1$ to tier $n$:
  \[
    \mathbb{E}[T_{n \to n+1}] = \PHI \cdot \mathbb{E}[T_{n-1 \to n}].
  \]
  Higher tiers require geometrically more effort — which is exactly right.
  The sovereign tier is genuinely rare, not merely labelled.
\end{corollary}

\subsection{Compound Multiplier}

At \tier{Sovereign} tier and above, contributors earn a compound multiplier
$m_c > 1$ on all future reward accrual.  This multiplier is itself governed
by the $\PHI$-scaling: $m_c = 1 + \PHIINV^k$ for some $k$ depending on
sustained contribution duration.  The practical effect is that a
\tier{Sovereign}-tier contributor who continues to contribute receives
$\PHI$-compounding rewards — not as a bonus, but as a structural consequence
of the pressure-to-reward conversion at high tier.

% ─────────────────────────────────────────────────────────────────────────────
\section{Collective Coherence: Kuramoto Career Synchronisation}
\label{sec:kuramoto}

\subsection{Phase Coherence Bonus}

Each contributor has a Kuramoto phase $\Theta(c, t) \in [0, 2\pi)$ that
advances at a natural rate $\omega_c$ plus coupling to the mean field:
\begin{equation}
  \frac{d\Theta(c)}{dt} = \omega_c
    + K \cdot r \cdot \sin(\psi - \Theta(c)),
  \label{eq:kuramoto_career}
\end{equation}
where $r \cdot e^{i\psi} = (1/N) \sum_j e^{i\Theta(j)}$ is the Kuramoto
order parameter.

Contributors whose phase is close to the mean phase $\psi$ — i.e.\ who
contribute \emph{when the organisation is most active} — receive a coherence
bonus proportional to $\cos(\psi - \Theta(c))$:
\begin{equation}
  \text{bonus}(c) = w_k \cdot \delta \cdot \max(0, \cos(\psi - \Theta(c))).
  \label{eq:bonus}
\end{equation}

\subsection{Collective Career Emergence}

\begin{theorem}[Collective Career Emergence]
  When the contributor population's Kuramoto order parameter $r > 0.95$
  (the OMNIS threshold), the aggregate contribution rate exceeds the sum of
  individual rates by a factor of $r / r_0$, where $r_0$ is the baseline
  order parameter at $K = 0$ (uncoupled).
\end{theorem}

\begin{proof}[Sketch]
  At $r > 0.95$, the organisation's effective coupling $K_\text{eff}$ is at
  its maximum (all contributors are phase-aligned).  The phase-coherence bonus
  $\cos(\psi - \Theta(c)) \approx 1$ for all contributors.  Therefore, every
  contribution earns its full bonus, and the total reward pool is
  $(1 + \text{bonus factor}) \times$ the uncoupled total. \qed
\end{proof}

\subsection{Interpretation}

This result provides a formal justification for an empirically observed
phenomenon: organisations where members work in tight temporal coordination
(sprints, stand-ups, synchronised releases) produce more than the sum of
their individually measured outputs.  The Kuramoto model explains this as
phase-coherence bonus: coordinated action earns a multiplier that
asynchronous action does not.

% ─────────────────────────────────────────────────────────────────────────────
\section{AI Agent Contributors}
\label{sec:agents}

\subsection{Why the Framework Is Agent-Type Agnostic}

The career flow model makes no assumption about the cognitive substrate of a
contributor.  It requires only: (i) a unique principal ID; (ii) a verifiable
contribution record; (iii) a Kuramoto phase.  An AI agent satisfies all three
as naturally as a human contributor:

\begin{itemize}[leftmargin=*]
  \item The agent's canister ID is its principal ID.
  \item All its calls are recorded on the ICP blockchain — immutably and
    verifiably.
  \item Its tick cadence (the frequency at which it executes \texttt{solverTick}
    or its equivalent) defines its natural Kuramoto frequency $\omega_c$.
\end{itemize}

\subsection{Ring Affinity and Agent-Specific Rewards}

AI agents deployed at different ring levels in the sovereign architecture
receive ring-affinity bonuses:
\[
  \text{ring bonus}(c) = w_k \cdot \delta \cdot \PHI^{12 - r_c},
\]
where $r_c \in \{1, \ldots, 12\}$ is the agent's ring (1 = sovereign core,
12 = public surface).  Agents at deeper rings (smaller $r_c$) are rarer, more
powerful, and earn exponentially higher rewards for equal output.

\subsection{Founder Agents}

Agents deployed by the founder at genesis receive a permanent founder bonus:
\[
  W(c)_\text{genesis} = \PHI \cdot S_0 = \PHI \approx 1.618
\]
instead of the standard $S_0 = 1.0$.  This bonus is locked at genesis and
cannot be removed.  It reflects the structural fact that an agent present at
genesis bootstrapped the system's initial coherence — a contribution that
cannot be replicated by later entrants.

% ─────────────────────────────────────────────────────────────────────────────
\section{Incentive Structure Analysis}
\label{sec:incentives}

\subsection{Contribution Kind Efficiency}

Given a fixed effort budget, a contributor maximises pressure accumulation
by prioritising $\contrib{Govern}$ contributions ($w_k = \PHI^3$) over
$\contrib{Gate}$ contributions ($w_k = \PHIINV$) — by a factor of
$\PHI^3 / \PHIINV = \PHI^4 \approx 6.85$.  This is intentional:
governance contributions are weighted highest because they shape the rules
governing everything else.  A contributor who governs well does more
structural good than a contributor who merely calls tools.

\subsection{The No-Drop Law as Risk Insurance}

Traditional career models subject contributors to reputation risk: a public
failure can permanently reduce standing.  Under the No-Drop Law, no single
event can permanently damage a contributor's career.  This reduces the
variance of career outcomes and encourages larger, riskier contributions:
contributors who know their floor is protected are more willing to attempt
high-variance high-reward contributions.

\begin{proposition}[Risk Encouragement Effect]
  Under the No-Drop Law, the optimal contribution size $\delta^*$ is weakly
  larger than in a model without the floor guarantee.
\end{proposition}

\begin{proof}[Sketch]
  Let $\delta^*_\text{no-floor}$ maximise expected utility in a model where
  failure reduces $W$.  Let $\delta^*_\text{floor}$ be the optimum under
  No-Drop.  Because failure has no downside under No-Drop (it does not reduce
  $W$, only fails to increase it), the expected cost of a large contribution
  is weakly lower.  Therefore $\delta^*_\text{floor} \ge \delta^*_\text{no-floor}$.
  \qed
\end{proof}

\subsection{Phase Coherence as Career Coordination Mechanism}

The Kuramoto coupling in~\eqref{eq:kuramoto_career} acts as an endogenous
coordination signal: contributors who observe that the order parameter is high
($r \approx 1$) face a strong incentive to align their phase with the mean
field, because the coherence bonus~\eqref{eq:bonus} is maximised at $\cos 0 = 1$.
This creates a self-reinforcing coordination cycle: high $r$ attracts coherent
contributors; coherent contributors increase $r$.

% ─────────────────────────────────────────────────────────────────────────────
\section{Discussion}
\label{sec:discussion}

\subsection{Relationship to Existing Models}

The career flow model synthesises three established frameworks:
(i) Hebbian learning (reputation non-decrease) from neuroscience~\cite{hebb1949};
(ii) Kuramoto synchronisation (collective coherence bonus) from
physics~\cite{kuramoto1984};
(iii) tournament theory (tier-based rewards) from labour
economics~\cite{lazear1981}.  The synthesis is novel in two respects: the
No-Drop constraint is stronger than the ``reputation does not decay'' assumption
in most reputational equilibrium models~\cite{diamond1989}; and the Kuramoto
coupling links individual career flows to collective organisational emergence
in a way that has not been previously formalised.

\subsection{Limitations}

\paragraph{Strategic contribution timing.}
Sophisticated contributors may observe the mean phase $\psi$ and time their
contributions to maximise the coherence bonus.  This is not necessarily
harmful — it aligns contributor timing with organisational rhythm — but it
may disadvantage contributors who cannot observe $\psi$ (e.g.\ contributors
in low-connectivity environments).

\paragraph{Contribution magnitude verification.}
The model assumes $\delta$ is verifiable (i.e.\ on-chain).  Contributions
whose magnitude requires off-chain verification (e.g.\ the quality of a
research paper) require an oracle mechanism not covered by this model.

\paragraph{Sybil resistance.}
The model assumes each contributor has one principal ID.  A contributor with
multiple IDs could split contributions to game tier thresholds or phase
bonuses.  Standard on-chain identity mechanisms (staking, NNS neuron
verification) address this but are not formalised here.

\subsection{Future Work}

\begin{itemize}[leftmargin=*]
  \item Formal analysis of the steady-state tier distribution under mixed
    human/AI contributor populations with heterogeneous natural frequencies.
  \item Empirical validation from the \textsc{Nova} organism's live
    contributor record after 12 months of operation.
  \item Extension to \emph{multi-organism} settings where contributors can
    transfer their career record between sovereign organisms — portability of
    the No-Drop guarantee across systems.
\end{itemize}

% ─────────────────────────────────────────────────────────────────────────────
\section{Conclusion}
\label{sec:conclusion}

A persistent AI organisation needs a career model that is objective, permanent,
no-drop, scalable, and equitable across agent types.  We have shown that the
combination of the No-Drop Law (Hebbian weight floor at $S_0 = 1.0$),
$\PHI$-scaled sovereign tier progression (pressure accumulates geometrically),
and Kuramoto career synchronisation (phase-coherent contributors earn a
collective bonus) provides exactly this model.

The No-Drop Law is a structural invariant, not a policy commitment: the
mathematics enforces it.  Tier progression is genuinely meritocratic: the
$\PHI$ spacing between tier thresholds ensures that higher tiers are
geometrically harder to achieve.  And the Kuramoto coupling ensures that
individual career success and collective organisational coherence are
co-optimised, not in tension.

In this model, a career is not a story told by managers about an employee.
It is a verifiable mathematical trajectory, accumulating on-chain, owned by
the contributor, irreversible, and open to human and AI agents on equal terms.

% ─────────────────────────────────────────────────────────────────────────────
\begin{thebibliography}{99}

\bibitem{hebb1949}
D.~O.~Hebb,
\emph{The Organization of Behavior},
Wiley, 1949.

\bibitem{kuramoto1984}
Y.~Kuramoto,
\emph{Chemical Oscillations, Waves, and Turbulence},
Springer-Verlag, 1984.

\bibitem{lazear1981}
E.~P.~Lazear, S.~Rosen,
``Rank-order tournaments as optimum labor contracts,''
\textit{J.\ Polit.\ Econ.}, vol.~89, no.~5, pp.~841--864, 1981.

\bibitem{diamond1989}
P.~Diamond,
``Reputation and politics: Career concerns of politicians under voter
uncertainty,''
\textit{Am.\ J.\ Polit.\ Sci.}, vol.~39, no.~1, pp.~161--198, 1995.

\bibitem{strogatz2000}
S.~H.~Strogatz,
``From Kuramoto to Crawford,''
\textit{Physica D}, vol.~143, pp.~1--20, 2000.

\bibitem{shapiro1983}
C.~Shapiro, J.~Stiglitz,
``Equilibrium unemployment as a worker discipline device,''
\textit{Am.\ Econ.\ Rev.}, vol.~74, no.~3, pp.~433--444, 1984.

\bibitem{spence1973}
M.~Spence,
``Job market signaling,''
\textit{Q.\ J.\ Econ.}, vol.~87, no.~3, pp.~355--374, 1973.

\bibitem{dfinity2022}
\text{DFINITY Foundation},
``The Internet Computer,''
Technical White Paper, 2022.

\end{thebibliography}

\end{document}
